\section{Materialized Task Example for Loop Engineering Evaluation}\label{app:example-task}

This appendix presents one materialized task and the loop trace metrics that \lhb{} records during evaluation, making the schema in Section~\ref{sec:design} and Appendix~\ref{app:task-schema} concrete. The example is a PR Sequences task drawn from an actively maintained open source repository. It connects materialization steps to loop trace metrics, including commit history snapshotting, per pull request diff extraction by role, dependency routing, ready frontier test release, and regression obligation retention. Under the contamination protocol of Appendix~\ref{app:contamination}, identifiable tokens are redacted and marked as \texttt{<repo>}, \texttt{<unit-id>}, and \texttt{<author>}.

\paragraph{Native artifact.}

The task originates from a sequence of \textbf{9} merged pull requests against \texttt{<repo>} over \textbf{4.2 months}, touching \textbf{37} source files and \textbf{14} test files. The base codebase at the entry commit comprises \textbf{18.4k} LoC across \textbf{142} files. Each pull request provides a description, a per file diff, and a CI run record from the upstream test suite, giving source evidence for implementation units, prerequisite relations, and executable obligations.

\paragraph{Preprocessing output.}

After preprocessing, the nine pull requests decompose into \textbf{12} atomic candidate units. Three original PRs are split when their diffs touch disjoint feature surfaces, and two adjacent PRs are merged when the second repairs a regression introduced by the first and is therefore dependent on that PR. The materialization extractor produces the aggregated segments under the prescribed extraction procedure, and the segments pass the fixed acceptance standard on the first attempt.

\paragraph{Materialized task bundle.}
The candidate is rewritten into the uniform schema of Appendix~\ref{app:task-schema}. The resulting on disk layout is:
\begin{promptbox}{Task bundle layout}
<repo>-<unit-id>/
  Dockerfile
  docker-compose.yaml
  run-tests.sh
  solution.sh
  gold-patch.diff
  task.yaml
  tests/
    unit_001/...
    unit_002/...
    ...
    unit_012/...
\end{promptbox}

The \texttt{Dockerfile} pins Python 3.11.6, three system packages absent from the upstream README but recovered from CI configuration, and a lockfile derived from the entry commit. The materialization stage verifies \texttt{gold-patch.diff} successfully after a single build attempt.

\paragraph{Per unit requirement.}
A representative unit's \texttt{task.yaml} entry (paraphrased and redacted) reads:
\begin{promptbox}{Per unit requirement (excerpt)}
unit_id: <unit-id>
requirement: |
  Extend the streaming response handler so that
  partial chunks emitted before a terminating
  sentinel are buffered into a single payload
  delivered to downstream consumers, preserving
  ordering across concurrent producers.
file_scope:
  - src/stream/handler.py
  - src/stream/buffer.py
symbol_scope:
  - StreamHandler.consume
  - StreamHandler._flush
acceptance:
  - concurrent producers preserve emission order
  - sentinel terminates buffering atomically
  - downstream observers see one payload per
    logical message
parents: [unit_003, unit_005]
tests:
  - tests/unit_007/test_ordering.py
  - tests/unit_007/test_sentinel.py
  - tests/unit_007/test_concurrent.py
\end{promptbox}

The requirement is phrased as an acceptance contract. Procedural directives, command level instructions, and references to the gold patch are omitted. The resulting loop trace records whether the evaluated loop infers ordering and integration constraints from the dependency graph and released tests without access to the hidden implementation path.

\paragraph{Recovered DAG.}

The recovered DAG encodes the routing and obligation structure recorded by the evaluation runtime. DAG recovery (Section~\ref{sec:dag-framework}) admits \textbf{18} edges over the 12 units under the unambiguous evidence rules. The longest dependency chain has depth \textbf{6}. The maximum out degree is \textbf{4}, corresponding to a foundational refactor unit imported by later units. Three units form an independent branch with zero incoming edges from the rest of the task, creating a parallel ready frontier midway through the PR sequence.

\paragraph{Trial verification.}
The materialized test suite passes the solvability, non-triviality, and discriminativeness trials of Section~\ref{sec:description-refine} on the first run for \textbf{10} of \textbf{12} units. Two units undergo one repair round during materialization. In one case a flaky timing assertion is replaced with a deterministic state check. In the other, a discriminativeness trial reveals that two predecessor tests already cover the unit's full contract, and the materialization pipeline strengthens the unit suite with an additional assertion targeting the new contract.

\paragraph{Evaluation trace excerpt.}

Under the protocol of Section~\ref{sec:protocol}, the recorded loop trace contains \textbf{31} checkpoints over the evaluation episode. The visitation sequence $\sigma$ differs from the human order $\sigma^*$ on \textbf{3} of 12 units, all within branches with partially ordered alternatives under the DAG. Two regressions are recorded against previously completed units, and one is recovered before the episode ends. The full test state matrix $S$ is omitted for brevity, while Appendix~\ref{app:metrics} defines the corresponding loop trace metrics.
\FloatBarrier
